\section{A precise route to a Lean formalization}
\label{sec:lean-plan}

\subsection{The proof-theoretic endpoint}

There are two different formalization goals.  The first is to prove the
numerical estimates in Lean's ambient mathematics.  The second, required
here, is to exhibit a derivation using only the axioms and inference rules
of intuitionistic primitive recursive arithmetic.  The first goal does not
imply the second: even a Lean proof containing no use of excluded middle
may use induction considerably stronger than quantifier-free induction.

The supplied ProveIt repository already provides a suitable deep embedding.
The following statements refer to commit
\texttt{b03925cccc76d55b9dc2a00864a5e0544b64aa7f}, inspected through the
repository tools.  The source directory is
\begin{quote}\small\ttfamily
Logic/PrimitiveRecursiveArithmetic/ConciseIPRA/\\
Lean/ConciseIPRA/
\end{quote}
Within it, \texttt{PR.lean} defines primitive-recursive function codes
\texttt{PR k}, scoped arithmetic terms \texttt{PTm n}, and the axiom family
\texttt{PRAAx}.  The function codes are generated by zero, successor,
projections, composition, and primitive recursion.  The induction axiom
requires a proof that its matrix is quantifier-free; the useful derived
rule is \texttt{PRA.qfInd}.  The module \texttt{Deduction.lean} defines
\texttt{Prf Ax Gamma phi}, intuitionistic first-order natural deduction
with equality.  Its constructors provide no excluded-middle rule.

The final theorem should therefore have the form
\begin{quote}\ttfamily
theorem pnt\_iPRA : Prf PRAAx [] pntSentence := ...
\end{quote}
Here \texttt{pntSentence} must be a \emph{defined closed formula}, with
defined primitive-recursive codes for its numerical functions.  It is not
enough to leave those codes as parameters, to prove merely that suitable
functions are primitive recursive, or to establish an asymptotic statement
about \texttt{Real.log} in Mathlib.

The existing module \texttt{NativeBiInterpretation.lean} supplies
\texttt{pra\_iff\_native}:
\[
 (\Gamma\vdash_{\mathrm{iPRA}}\varphi)
 \quad\Longleftrightarrow\quad
 \operatorname{NPrf}
   (\Gamma.\operatorname{map}\operatorname{NFm.ofPRA})
   (\operatorname{NFm.ofPRA}(\varphi)).
\]
Consequently, a completed arithmetic derivation transfers directly to the
repository's concise native calculus.  This reuses the foundational
equivalence without treating it as a source of new number-theoretic
estimates. The inspected Git blobs are recorded explicitly for reproduction:
\begin{quote}\small
\texttt{PR.lean}:\\
\texttt{3251631592070d9ed3687cc0bf7dbd5eb452e46f}\\[3pt]
\texttt{NativeBiInterpretation.lean}:\\
\texttt{ccfe499b9687ea0e176f832686d8e492fe15544b}.
\end{quote}

\subsection{An exact arithmetic statement}

For an integer $n\geq 1$, put
\[
 F(n)=(n^2)!,\qquad
 A(n)=\sum_{j=n+1}^{n^2}\frac{F(n)}{j},\qquad
 \ell(n)=\frac{A(n)}{F(n)}.
\]
Every quotient in $A(n)$ is exact.  Thus both $A$ and $F$ are functions on
natural numbers, and $\ell(n)=H_{n^2}-H_n$.  The main decidable matrix is
\[
 k\,\operatorname{dist}\bigl(\pi(n)A(n),nF(n)\bigr)\leq nF(n),
 \qquad
 \operatorname{dist}(a,b)=(a\mathbin{\dotminus}b)+(b\mathbin{\dotminus}a),
\]
where $\dotminus$ is truncated subtraction.  For $k>0$ and $n\geq 2$ this is
exactly
\[
 \left|\frac{\pi(n)\ell(n)}n-1\right|\leq\frac1k.
\]
Use the explicit modulus
\[
 N(0)=2,\qquad
 q=8k,\quad d=2^{20}q^3,\quad
 N(k)=2^{\,2^{d^2+2}}\quad(k>0).
\]
The Boolean checker returns true if $k=0$, if $n<N(k)$, or if $n<2$;
otherwise it checks the displayed integer inequality.  The closed target
is simply the universal closure of ``this checker returns true.''  It
contains no real-number objects or quantification over functions.

All its numerical constructions admit direct primitive-recursive codes.
Primality is a bounded test over $d\leq n$; $\pi$ is a running sum of its
characteristic function; factorial, products, and sums are primitive
recursions.  Division and truncated subtraction are bounded-search
constructions.  Exponentiation is primitive recursive, so the displayed
fixed nesting of exponentials causes no foundational difficulty.  Large
values affect computational practicality, not the permitted recursion
scheme.  In particular, the formalization should never attempt to expand
the closed numeral $N(1)$ during evaluation.

\subsection{Finite ceilings and arithmetic encodings}

The quantitative argument iterates an error bound.  Its implementation
must not apply induction directly to
\[
 \forall n\geq P_j\quad |\psi(n)-n|\leq a_j n,
\]
because this is not a quantifier-free induction matrix.  Instead fix an
arbitrary finite ceiling $X$ before starting the iteration.  Encode the
assertions for every relevant $n\leq X$ by a single Boolean conjunction.
For any fixed primitive-recursive Boolean predicate $b$, this conjunction
is generated by
\[
 B(0)=b(0),\qquad B(t+1)=B(t)\wedge b(t+1).
\]
Its defining equations and correctness are established by quantifier-free
induction.  The point is that $b$ is compiled as a \emph{fixed arithmetic
code}; the object language does not quantify over an arbitrary function
parameter.

For the piecewise-constant Chebyshev function, a convenient finite
invariant records both endpoint estimates.  If $U_\psi(n)$ and
$L_\psi(n)$ are rational upper and lower enclosures for $\psi(n)$, require
\[
 U_\psi(n)\leq(1+a_j)n,
 \qquad
 L_\psi(n)\geq(1-a_j)(n+1)
\]
for the relevant integers $n\leq X$.  These inequalities control the
normalized error throughout each interval $[n,n+1)$, including its left
limit at the right endpoint.  All denominators are cleared before the
Boolean test is formed.  The contraction step uses only arguments below
its current finite ceiling; its rational mesh, bounded selection of
indices, and numerical precision are also given primitive-recursive
bounds.  One then applies \texttt{PRA.qfInd} to this Boolean invariant,
and finally substitutes $X=n$.  This substitution supplies the desired
unbounded universal conclusion without using a stronger induction schema.

Finite rational enclosures are part of the proof, rather than an appeal
to numerical experiments.  Each analytic abbreviation must come with a
proved finite error estimate.  Rational sign decisions use separated
thresholds, so equality of an unknown real number to a threshold is never
tested.  The surplus in the contraction estimate pays for the enclosure
errors and the passage to endpoint bounds.  A useful implementation order
is to prove these interval contracts first, then to reuse them in every
finite summation and mesh argument.

\subsection{Module milestones and verification}

The following modules are prospective divisions of the remaining work;
their names do not claim that such completed modules already exist.
\begin{enumerate}
\item \textbf{Arithmetic codes and their defining equations.}
Compile the exact checker, including $N,A,F,\pi$, into a concrete term of
\texttt{PR 2}.  Prove that its ordinary interpretation agrees with the
executable checker.  Build reusable object-level proofs for bounded sums,
products, searches, and Boolean conjunctions.

\item \textbf{Rational intervals.}
Develop finite logarithm enclosures, exact rational arithmetic, and the
harmonic estimate for $\ell(n)$.  Every precision parameter has a specified
primitive-recursive truncation bound.  Separate these verified interval
contracts from the later conventional interpretation in the real numbers.

\item \textbf{Finite number-theoretic identities.}
Formalize factorization, the identities for the von Mangoldt and M\"obius
functions, the elementary Chebyshev bounds, and the finite summation
identities used in Selberg's formula.  Prove their rational error versions
with the constants used in the human-readable argument.

\item \textbf{The quantitative contraction.}
Formalize the bounded geometric mesh, the separated sign alternatives,
the deficit estimate, and the threshold recurrence.  Carry a finite
ceiling $X$ throughout.  The crucial outputs are explicit transformations
of the endpoint Boolean invariants, not an invocation of an unbounded
limsup argument.

\item \textbf{The modulus and the counting function.}
Verify the finite number of contraction steps and the majorization by
$N(k)$.  Complete the transfers from $\psi$ to $\vartheta$ and to $\pi$,
then replace the logarithm by $A(n)/F(n)$ with the stated error budget.

\item \textbf{The object-level certificate.}
Assemble the previous derivations into
\texttt{Prf PRAAx [] pntSentence}.  Transfer this certificate through
\texttt{pra\_iff\_native}.  Reject any remaining proof hole, additional
arithmetic axiom, or strengthened induction rule.  Record
\texttt{\#print axioms} for the exported certificate and audit the
definition of the final sentence.

\item \textbf{The semantic PNT corollary.}
Only after the arithmetic result is established, recover the conventional
statement using standard natural-number semantics and the elementary
comparison with the real logarithm.  This last theorem may be presented
using Mathlib's usual real-analysis interface; its logical status is
distinct from the object-level certificate.
\end{enumerate}

The role of the final semantic step is particularly transparent.  The
elementary inequalities for the integral of $1/t$ give
\[
 0\leq\log n-\ell(n)\leq\frac1n.
\]
Since $\pi(n)\leq n$, replacing $\ell(n)$ by $\log n$ changes the normalized
quantity by at most $1/n$.  Applying the arithmetic theorem at precision
$2k$ and taking $n\geq 2k$ yields error at most $1/k$ for
$\pi(n)\log n/n$.  Thus, for example,
\[
 \max\{N(2k),2k,2\}
\]
is a conventional modulus obtained directly from the exact arithmetic
sentence.  No hidden choice of a sufficiently large threshold occurs.

\subsection{Status of the supplied Lean artifacts}

\texttt{lean/RationalPNT.lean} contains executable definitions of the
natural-number checker and the explicit modulus, together with small
generic proofs about bounded Boolean conjunctions, prime-counting bounds,
positive factorials, and the equivalence of the finite-ceiling
specification.  It imports only \texttt{Std}.  The definition
\texttt{ArithmeticPNT} states the target; it is not a theorem proving it.

\texttt{lean/ObjectTarget.lean} supplies a precise constructor for the
closed object-language sentence and generic wrappers around the
repository's native/iPRA proof transfer.  The arithmetic code remains a
parameter in these wrappers until the compilation milestone is completed.
No PNT axiom is introduced, and no wrapper is presented as the missing
number-theoretic derivation.

\textbf{These Lean files are an outline, not a completed formalization.}
They were not compiled in the preparation environment, where no Lean,
Lake, or Elan executable was found.  The mathematical estimates and their
object-level derivations remain the formalization work described above.
The accompanying \texttt{lean/STATUS.txt} records this limitation.  A
successful future build must establish both a faithful arithmetic statement
and a derivation in the specified object theory; the absence of
\texttt{sorry} in a collection of definitions alone would not establish
either main theorem.
